\section{Introduction}

The Islamic financial system is logically distinct from the conventional financial system because its motivation is grounded in Islamic Shariah principles. Although these principles are normatively stable, their practical application needs to be reconsidered in light of contemporary economic conditions and the technological infrastructures that increasingly shape financial markets. In this context, Islamic finance should not only preserve its ethical and legal foundations, but also engage with modern analytical tools that can improve investment decision-making, risk management and portfolio construction.

This argument is closely related to the concept of the Islamic Moral Economy. Rather than reducing Shariah compliance to a narrow screening mechanism, the Islamic Moral Economy emphasises ethical purpose, social welfare and human-centred development \citep{asutay2007conceptualisation, asutay2012conceptualising}. From this perspective, the use of machine learning (ML) and deep learning (DL) in Islamic finance can be understood as more than a technical exercise. These methods can support the development of investment frameworks that are both computationally advanced and ethically aligned.

Portfolio optimisation is traditionally rooted in the mean-variance framework developed by \citet{markowitz1952portfolio}. Its main objective is to maximise expected return for a given level of risk, or equivalently to minimise risk for a given level of expected return. Although this framework emerged from conventional financial theory, it can be adapted to Shariah-compliant investing when the investment universe is restricted to halal equities that satisfy Shariah screening criteria. Therefore, portfolio optimisation operates within an ethical boundary that limits exposure to prohibited sectors and financial structures.

Nevertheless, Islamic finance has been criticised for imitating conventional financial models without sufficiently realising its broader social and ethical objectives \citep{asutay2007conceptualisation}. This paper responds to that concern by examining whether advanced ML and DL architectures can strengthen the operational capacity of Shariah-compliant investment management. In particular, predictive models can improve return estimation, while sparse optimisation techniques can help identify manageable portfolios within a high-dimensional market environment. In this way, computational finance can support the practical implementation of Islamic investment principles.

To operationalise this approach, the paper provides an empirical analysis of the BIST Participation All Index (BIST XKTUM) over the period from 2007 Q1 to 2022 Q4. The study uses daily closing prices of 163 Shariah-compliant stocks. The first stage of the methodology applies six predictive models: Long Short-Term Memory (LSTM), Bidirectional LSTM (Bi-LSTM), Gated Recurrent Unit (GRU), Extreme Gradient Boosting (XGBoost), Random Forest (RF) and Support Vector Machine (SVM). These models are selected to compare sequence-based deep learning models with non-sequential machine learning models.

The second stage integrates the predicted returns into a mean-variance optimisation framework. L1 regularisation is applied to support sparse asset selection and reduce portfolio complexity. This is particularly relevant in large Shariah-compliant equity universes, where investors may need to construct a manageable portfolio while maintaining risk-return efficiency. Due to the computational cost of hyperparameter tuning across 163 stocks, NVIDIA V100 GPUs are used to enable efficient model training and grid-search optimisation. In this paper, GPU acceleration is treated as a computational support mechanism rather than a direct cause of predictive accuracy.

The contribution of this paper is threefold. First, it compares ML and DL models in the context of Shariah-compliant stock price prediction. Second, it links predictive modelling with portfolio optimisation rather than evaluating forecasting accuracy alone. Third, it introduces a sparse optimisation approach for managing large-scale Shariah-compliant portfolios. These contributions address a gap in the literature, where Islamic fintech and Islamic investment studies often remain conceptual, while empirical studies combining predictive analytics, hyperparameter optimisation and Shariah-compliant portfolio construction remain limited.

The rest of the paper is structured as follows. Section 2 reviews the literature on Islamic Moral Economy, portfolio optimisation and machine learning-based stock prediction. Section 3 presents the data and methodology. Section 4 discusses the comparative prediction results. Section 5 presents the portfolio optimisation results. Section 6 concludes with limitations, practical implications and future research directions.